% block: multigrid-heading page:31 type:section
\subsection{\BT{Multigrid Methods}{多重网格方法}}\label{sec:multigrid-methods}
% block: multigrid-principle page:31 type:ordinary-paragraph
\BT{Multigrid methods fall into a class of iterative methods that have been developed for solving linear systems arising from elliptic partial differential equations. Multigrid methods find a good approximate solution of the linear systems from the information of some rough approximate solutions of smaller linear systems (arising from hierarchical coarser meshes) and an accurate approximate solution of the smallest linear systems (arising from the coarsest mesh). For obtaining the rough approximate solutions, stationary iterative methods (especially the Gauss--Seidel method) are often used, which are called smoothers. The name smoother comes from the fact that the Gauss--Seidel method damps quickly the high frequency components of the error or residuals for the linear systems, and the multigrid methods make the most of this property.}{多重网格方法是一类为求解椭圆型偏微分方程所产生的线性方程组而发展起来的迭代方法. 它利用较小线性方程组 (来自分层的较粗网格) 的若干粗略近似解, 以及最小线性方程组 (来自最粗网格) 的精确近似解所提供的信息, 求得原线性方程组的良好近似解. 为得到粗略近似解, 常使用称为光滑子的定常迭代方法 (尤其是 Gauss--Seidel 方法). 光滑子这一名称来自这样一个事实: Gauss--Seidel 方法能够快速衰减线性方程组的误差或残差中的高频分量, 而多重网格方法充分利用了这一性质.}
% block: multigrid-convergence page:31 type:ordinary-paragraph
\BT{The notable feature of the multigrid methods is that the convergence rate does not depend on the mesh size, though stationary iterative methods and Krylov subspace methods do. The multigrid methods are later extended to solving other linear systems arising from not only elliptic ones but other models.}{多重网格方法的显著特点是收敛率不依赖于网格尺寸, 而定常迭代方法和 Krylov 子空间方法的收敛率依赖于网格尺寸. 后来, 多重网格方法被推广到求解其他线性方程组, 不仅包括椭圆型问题产生的方程组, 也包括其他模型产生的方程组.}
% block: multigrid-books page:31 type:ordinary-paragraph
\BT{Multigrid methods are the subject of active research, and there are many excellent books on multigrid methods, some of them by:}{多重网格方法是一个研究活跃的课题, 已有许多优秀著作, 其中一些作者为:}
\begin{itemize}
\item Hackbusch \BookCite{91},
\item Wesseling \BookCite{204},
\item Bramble \BookCite{27},
\item Briggs et al. \BookCite{28},
\item Trottenberg et al. \BookCite{193},
\item Lottes \BookCite{125}.
\end{itemize}
\BT{Wolfgang Hackbusch's book presents not only the basic concepts of multigrid methods using simple one-dimensional differential equations but also describes algorithms, their program codes, and detailed convergence analysis. James Lottes's book focuses mainly on algebraic multigrid methods.}{Wolfgang Hackbusch 的著作不仅借助简单的一维微分方程介绍多重网格方法的基本概念, 还给出算法, 程序代码以及详细的收敛分析. James Lottes 的著作则主要讨论代数多重网格方法.}
